\documentclass[12pt]{article}
\usepackage[utf8]{inputenc}
\usepackage[T1]{fontenc}
\usepackage{hyperref}
\title{Morerebi Language Resources for K-12 Curriculum Development}
\author{Compiled by Claude Sonnet 5.5 (claude-sonnet-5-5) for Andr\'as Kornai}
\date{}
\begin{document}
\maketitle

\section{Introduction}
Morerebi (ISO 639-3 xmo; Glottocode more1262) is an extremely poorly documented language of Brazil. Joshua Project and Global Recordings Network give a figure of about 100 people, in Amazonas state (Marmelos and Rio Preto) and Mato Grosso, citing Ethnologue. Ethnologue places it in the Tupian family, while Glottolog lists it as unattested and describes it as an uncontacted group of the Rio Preto and Marmelos area. One source suggests a possible link with Tenharim. Joshua Project says the group may have assimilated into surrounding groups and that no written materials or Bible translation exist.

\section{Primary Educational Resources}
None found. No primer, textbook or school programme exists as far as these searches show.

\section{Linguistic and Reference Materials}
No grammar, dictionary or word list turned up. Only catalogue entries were found:
\begin{itemize}
\item Glottolog entry: \url{https://glottolog.org/resource/languoid/iso/xmo}
\end{itemize}

\section{Religious and Community Resources}
\begin{itemize}
\item Joshua Project people-group profile (population and status only): \url{https://joshuaproject.net/people_groups/13811}
\end{itemize}

\section{Digital and Audio Resources}
\begin{itemize}
\item Global Recordings Network language page (the search result showed only a language-selection page; I did not verify whether any recordings exist): \url{https://globalrecordings.net/en/language/xmo}
\end{itemize}

\section{Curriculum Development Context}
There is no basis for a curriculum at present. Nothing documents the language, its speakers' current status is uncertain, and the group may be uncontacted. Any work would have to respect the wishes of the people concerned and involve the Brazilian indigenous authorities, and the first step would be basic documentation, if the community wants it. I found no evidence of community or institutional activity.

\end{document}